% !TEX root = ../../main.tex

\section{研究背景}

随着移动互联网、智能终端和人工智能技术的发展，面向特殊群体的信息无障碍服务逐渐从传统的信息展示与人工辅助，发展为融合移动应用、智能识别、多媒体资源和个性化反馈的综合服务形态。
《中华人民共和国无障碍环境建设法》明确提出，应当为残疾人、老年人自主安全地通行道路、获取和交流信息以及获得社会服务等提供便利，并对无障碍信息交流和无障碍社会服务提出了制度要求\cite{accessibility-law-2023}。
听障人士在日常交流、教育学习和社会参与过程中，常常面临语音信息获取受限、口语表达训练困难、手语学习资源分散以及与健听人群沟通成本较高等问题。
手语作为听障群体重要的视觉语言形式，具有自然、直观和表达丰富等特点；
口语训练、手势训练和视觉化动作示范也能够在一定程度上帮助用户提升交流能力和训练效率。
因此，构建面向听障人士的数字化交互训练系统，具有较强的现实意义和应用价值。

近年来，残疾人康复服务与信息无障碍建设持续推进。
中国残疾人联合会发布的《2024 年残疾人事业发展统计公报》显示，2024 年我国共有 866.8 万残疾人得到基本康复服务，其中接受基本康复服务的持证残疾人中包括 70.6 万听力残疾人\cite{cdpf-statistical-communique-2024}。
世界卫生组织相关资料也显示，全球范围内听力损失人群数量较大，听力损失会对交流、教育、就业和社会参与产生持续影响\cite{who-hearing-loss-2026}。
这些数据说明，听障群体的康复训练和辅助交流需求具有现实基础，利用移动终端和人工智能技术提升训练服务质量具有较好的应用前景。

传统听障训练方式较多依赖线下教学、人工示范和重复练习，训练过程容易受到时间、场地、师资和资源数量的限制。
对于普通用户而言，若缺少持续的标准动作示范和即时反馈，训练过程往往难以形成闭环；
对于系统维护人员而言，若缺少统一的资源管理、样本管理和模型版本管理能力，训练内容和识别模型也难以持续更新。
随着计算机视觉、关键点检测和深度学习技术的发展，基于摄像头视频的手势识别逐渐具备工程应用基础；
同时，基于移动端的训练应用能够降低用户使用门槛，使用户在更日常化的场景中完成课程学习、动作练习和结果查看。

HarmonyOS 作为面向多终端场景的新一代操作系统，为移动端训练应用的开发提供了较好的平台基础。
HarmonyOS 系统面向多终端提供“一次开发，多端部署”的能力，有助于开发者基于一套设计构建多端可运行的应用\cite{harmonyos-multi-device-deployment}。
其 ArkTS 与 ArkUI 技术体系支持声明式界面开发、状态驱动渲染和较自然的移动端交互组织，适合承载课程浏览、训练流程、相机采集、视频上传和结果反馈等功能。
另一方面，Spring Boot、Vue、MySQL、Redis、MinIO 等成熟工程技术能够支撑业务数据管理、后台资源维护和文件资源存储；
MediaPipe 与 TensorFlow 则为手部关键点提取、姿态分析和时序动作分类提供了可实现的技术路径\cite{lugaresi-mediapipe-2019,zhang-mediapipe-hands-2020}。

基于上述背景，本文设计并实现了基于 HarmonyOS 的听障人士交互训练系统 HearBridge。
系统以 HarmonyOS 移动端作为用户入口，结合 Spring Boot 服务端、Vue 后台管理端和 Python 手势识别服务，形成面向课程学习、手势训练、实时识别、句子视频识别、样本管理和模型发布的完整应用体系。
与单纯展示课程资源或单独进行模型实验不同，本文更强调训练应用、后台资源管理和识别服务之间的协同工作，使系统既能够面向用户提供较完整的训练体验，也能够面向后续模型优化和资源扩充提供工程支撑。
